\chapter{Introduction}
\label{cap:introduction}

% TODO: Expand with actual thesis context and motivation

\section{Context and motivation}
\label{sec:context}

Autonomous navigation in agricultural environments requires robots to follow crop rows reliably under variable field conditions. Vision-based guidance uses cameras to identify the navigable path between rows, but lighting changes, occlusions, weeds, and plant growth affect image appearance. Manual pixel-level annotation for training segmentation models is time-consuming and must cover multiple crops and conditions.

This thesis addresses these challenges through a combination of learning-based semantic segmentation and geometric reference extraction. It investigates model-assisted annotation to reduce manual labeling effort, compares segmentation architectures suitable for real-time robot perception, and develops geometric methods that convert corridor masks into local navigation references.

\section{Objectives of the thesis}
\label{sec:objectives}

The main objectives of this work are:
\begin{itemize}
\item Compare manual and SAM~3-assisted annotation strategies for agricultural segmentation datasets spanning lettuce, tomato, and cucumber.
\item Select a segmentation architecture balancing accuracy and inference speed for robotic deployment, through offline evaluation of FCN, DeepLabV3, and LR-ASPP on manually annotated lettuce images.
\item Evaluate the effect of output label configurations (four-class, three-class, and binary path) on path segmentation accuracy.
\item Implement and compare two edge-based geometric extractors that convert segmentation masks into local image-space navigation references.
\item Integrate the selected perception and geometric components on a robot platform for field validation.
\end{itemize}

\section{Organization of the work}
\label{sec:organization}

The thesis is organized as follows:

\textbf{Chapter~\ref{cap:background}} introduces the agricultural navigation problem, the role of vision-based local guidance, and the segmentation-to-control pipeline structure.

\textbf{Chapter~\ref{cap:stato_arte}} reviews vision-based methods for agricultural row navigation, covering classical and learning-based segmentation approaches, geometric extraction techniques, and experimental validation in related work.

\textbf{Chapter~\ref{cap:segmentation}} describes dataset preparation with manual and SAM~3-assisted annotation, the comparison of three segmentation architectures on lettuce, and the evaluation of supervision strategies with different output label configurations.

\textbf{Chapter~\ref{cap:geometric_reference}} presents the two geometric extractors that convert corridor masks into navigation references through horizontal-strip interval selection and independent edge fitting.

\textbf{Chapter~\ref{cap:platform}} describes the robot hardware, camera configuration, software integration using ROS, and the deployment of the trained segmentation model.

\textbf{Chapter~\ref{cap:experiments}} reports offline and field validation results for segmentation accuracy, geometric reference extraction, and closed-loop navigation performance.

\textbf{Chapter~\ref{cap:conclusions}} summarizes the main findings, discusses limitations, and identifies directions for future work.
